\scanpage{38}

2) Premières propriétés [concernant $\pi_0$, $\pi_1$, Pic]

\uline{Proposition} Supposons que
$H^0(X,O_X)\mathbin{-}H^0(X_{/Y},O_{X/Y})$
soit injectif [bijectif] alors
$\pi_0(Y)\longrightarrow\pi_0(X)$ est surjectif [bijectif].
En particulier $(LG)_{-1}$ [resp $(LG)_0$] implique que
$\pi_0(Y)\longrightarrow\pi_0(X)$ est surjectif [bijectif].

\uline{Proposition} Supposons que $(X,Y)$ satisfasse
$LG$ (resp. $LG$ stricte).
\inserted{et $\unsure{\eta|\xi}$ un pt géom. de $Y$
(\unsure{connexe})}
Alors $\pi_1(Y)\mathbin{-}\pi_1(X)$ est surjectif [resp.
\[
\pi_1(Y,\unsure{\eta|\xi})\longrightarrow
\mathop{\mathrm{lim}}_{\substack{\longleftarrow\\
U\text{ voisinage}\\\text{ouvert de }Y}}
\pi_1(U,\unsure{\eta|\xi})
\]
est un hom \deleted{surjectif [resp.} bijectif]].

Résulte des faits plus généraux

\uline{\deleted{Corollaire}} \inserted{Lemme}
Supposons $LG$. \deleted{+ Alors \unsure{pour} $\forall$}
$F$ loc libre sur $X$;
(i) tt structure d'alg.
\deleted{[\unsure{resp.} \unsure{associative}, à unité, resp. commutative]}
\deleted{\unsure{connue}} sur $\hat F$ provient d'une structure
d'algèbre unique sur $F$, qui est ass (resp unitaire,
resp. comm.) ssi celle de $\hat F$ l'est.

\deleted{\unsure{de même que}}
\deleted{(ii) les décompositions de $F$ \deleted{\unsure{sont aussi}}
\inserted{en} produits de \inserted{deux} algèbres non nulles
correspondent aux décompositions analogues de $\hat F$
[\unsure{en liaison avec|en biject. avec} \unsure{celles de} $\unsure{F(n)|\Gamma(X)}$].}
Donc si $X'=\operatorname{Spec}(F)$, alors $X'$ connexe
ssi $\hat X'_{/Y}$ connexe

\margin{(iii) Un hom \inserted{$F\longrightarrow G$}
\deleted{de préschémas} de Modules munis de structures
d'alg. est un hom d'alg. ssi
$\hat F_{/Y}\longrightarrow G_{/Y}$ l'est.}

\margin{\deleted{\uline{Lemme}
[\unsure{L'immersion}, \unsure{la}\illegible{1}
dans un \unsure{voisinage} \unsure{de}\illegible{1}
\unsure{particulier} \unsure{fini} \unsure{ét.}
\unsure{donne} un foncteur \uline{pl. fidèle}
de catégories \unsure{étales|d'actions}]}}

\margin{37}
